\documentclass[10pt,a4paper]{amsart}
\usepackage[margin=19mm]{geometry}
\usepackage{amsmath,amssymb,mathtools}
\usepackage[scr=rsfs,cal=euler]{mathalfa}
\usepackage{xfrac}
\usepackage[unicode,hidelinks,bookmarks=false]{hyperref}
\newcommand{\ie}{i.e.\ }
\newcommand{\cf}{cf.\ }
\newcommand{\resp}{respectively\ }
\newcommand{\Subsection}{\subsection}
\providecommand{\nobreakdash}{\nobreak\hbox{-}}
\setlength{\emergencystretch}{2em}
\multlinegap0pt
\usepackage{bm}
\usepackage{bbm}
\usepackage{dsfont}
\usepackage{xspace}
\usepackage{cancel}
\usepackage{mathtools}
\usepackage[shortlabels]{enumitem}
\usepackage{amsthm}
\makeatletter
\@ifundefined{theorem}{\newtheorem{theorem}{Theorem}}{}
\@ifundefined{claim}{\newtheorem{claim}[theorem]{Claim}}{}
\@ifundefined{lemma}{\newtheorem{lemma}[theorem]{Lemma}}{}
\@ifundefined{observation}{\newtheorem{observation}[theorem]{Observation}}{}
\makeatother
\newcommand{\id}{{0}}%{\mathrm{id}}
\title[On the Graham--Sloane harmonious labelling conjecture]{On the Graham--Sloane harmonious labelling conjecture}
\author{Alp M\"uyesser, Alexey Pokrovskiy}
\date{Source: arXiv:2509.05280v2 (CC BY 4.0)}
\begin{document}
\maketitle

\noindent\emph{Source and licence.} On the Graham--Sloane harmonious labelling conjecture. Version arXiv:2509.05280v2 (CC BY 4.0), distributed under the Creative Commons Attribution 4.0 International licence, as recorded in the arXiv licence metadata for this exact version and shown on the abstract page https://arxiv.org/abs/2509.05280v2. Full rights notice: licenses/arxiv-2509.05280-CC-BY-4.0.txt.\par\smallskip

\noindent\textit{Editorial scope.} Counted unit: the Theorem whose source label is \texttt{Theorem\_harmonious\_tree\_characterization} in the article (source0.tex lines 786--794, source label \texttt{Theorem\_harmonious\_tree\_characterization}), delivered together with the article's own complete original proof of it (source0.tex lines 795--866, one \texttt{proof} environment of 72 lines). No mathematics is written, repaired or re-derived: statement and proof are the article's own text line for line. Required context retained uncounted: Observation\_small\_core (lines 185--187); Theorem\_embed\_tree\_prescribed\_sets (lines 193--200); Lemma\_kastv\_simple\_tree (lines 729--734); Lemma\_y1y2\_separate (lines 748--750); Lemma\_y1y2...ys\_separate (lines 760--763); Lemma\_zero\_sum\_sets\_Z2k (lines 771--774). Every definition, display and label of those lines is retained unchanged. Omitted from this package and not redistributed: 1--184, 188--192, 201--728, 735--747, 751--759, 764--770, 775--785, 867--1172. Nothing inside a retained excerpt is omitted or altered: every retained body is delivered exactly as the article's lines are written, so no omission receipt of any kind accompanies this package. Citations: the retained excerpts carry no citation command at all, so no bibliography entry of the article is transcribed and no reference is redistributed. Reference adaptation: none was needed and none was made; every reference inside the delivered unit resolves inside it, so no unresolved reference or citation is produced. Printed slips kept literal: no printed word, symbol, hypothesis or number of the counted statement or of its proof was altered. Layout and class: the article's own class is \texttt{amsart} with packages including graphicx, geometry, hyperref, amsthm, thm-restate, amssymb, amsmath, cleveref, soul, comment, setspace, xcolor, enumitem, listings; the wrapper is the lane's pinned \texttt{amsart} editorial wrapper at 10pt on A4, which restates the theorem-like environments the retained text needs and adds the article's own macro definitions that the retained text needs (\texttt{\textbackslash id}), with extra wrapper packages (bm, bbm, dsfont, xspace, cancel, mathtools, enumitem). The delivered numbering is the unit's own. Assets: no retained span contains a figure environment, a graphics-inclusion command, a PDF-page inclusion, a table or a TikZ picture; no image file is copied into this package, the package declares no asset dependency and no image file is redistributed with it. Licence: version arXiv:2509.05280v2 of this article is distributed under the Creative Commons Attribution 4.0 International licence, as recorded in the arXiv licence metadata for this exact version and shown on the abstract page https://arxiv.org/abs/2509.05280v2. A full rights notice is delivered at licenses/arxiv-2509.05280-CC-BY-4.0.txt. Characters: every retained excerpt is pure ASCII, so the delivered engine, XeLaTeX with its default Latin Modern text font, reports no missing character.
\begin{observation}\label{Observation_small_core}
Every tree has a core $T_{core}$ of order $\leq 12\Delta(T)$.  If $\Delta(T)\leq \sqrt{|V(T)}/12$, then there is a core of order exactly $12\Delta(T)$.
\end{observation}

\begin{theorem}\label{Theorem_embed_tree_prescribed_sets}
Let $\Delta^{-1}\gg\mu\gg n^{-1}$.
Let $G$ be a group and $T$ a bounded degree tree with $\Delta(T)=\Delta$. Let $V_{target},C_{target}\subseteq G$ with $|T|=|V_{target}|=|C_{target}|+1\geq (1-n^{-\mu})n$. In the case $G=\mathbb{Z}_2^k$, assume $\id \not\in C_{target}$. Let $T_{core}$ be a core of $T$ of size $\leq n^{1-\mu}$. The following are equivalent.
\begin{enumerate}[(i)]
\item  $T$ has a rainbow  embedding $f$ into $(V_{target},C_{target})$.
\item There is a rainbow embedding $\phi$ of $T_{core}$ into $(V_{target},C_{target})$ with $\sum_{v\in V(T_{core})}d_T(v)\phi(v)=\sum C_{target}$ and $\sum \phi(V(T_{core}))=\sum V_{target}$.
\end{enumerate} 
\end{theorem}

\begin{lemma}\label{Lemma_kastv_simple_tree}
Let $\Delta^{-1}, k^{-1}\gg n^{-1}$, $G$ an order $n$ abelian group and $T$ a graph with $\Delta(T)\leq \Delta$. Let $v=(v_1, \dots, v_k)$ be a sequence of distinct vertices in $T$ having $d_T(v_i)\not\equiv 1 \pmod G$. Then there exists some simple $x\in G^{k}$ with $x \ast_T v$ simple unless:
\begin{itemize}
\item[$(\ast)$] $k=2$, $d_T(v_1), d_T(v_2)\equiv 0 \pmod G$, and $v_1v_2\in E(T)$.
\end{itemize}
\end{lemma}

\begin{lemma}\label{Lemma_y1y2_separate}
Let $G\neq \mathbb{Z}_2^k$ be an order $n$ abelian group. Then, for any $g\in G$, there are $> n/2$ solutions to $y_1+y_2=g$ with $y_1\neq y_2$. In particular, for any $F_1\subseteq G$ with $|F_1|< n/4$, there is such a solution with $y_1, y_2\not\in F_1$.
\end{lemma}

\begin{lemma}\label{Lemma_y1y2...ys_separate}
Let $C^{-1}, s^{-1}\gg n^{-1}$ with $s\geq 3$.
Let $G$ be an abelian group. Then, for any $g\in G$ and $F_1, F_2\subseteq G$ with $|F_1|, |F_2|\leq C$, there is a solution to $y_1+y_2+\dots+y_s=g$ with $y_i\not\in F_1, y_{i}-y_j\not\in F_2$ for all $i\neq j$.
\end{lemma}

\begin{lemma}\label{Lemma_zero_sum_sets_Z2k}
Let $T$ be a graph whose vertex set is partitioned $V(T)=A\cup B$. Suppose that $k\gg |A|+|B|\ge 10$, $|A|, |B|\ne 2$ and if $|A|=4$ or $|B|=4$ then $T[A]$ or $T[B]$ has no perfect matching (respectively). Then, there exists a rainbow embedding $\phi: T \to \mathbb{Z}_2^k$ also satisfying that $\sum \phi(A)=\sum \phi(B)=0$.
%Let  $a,b\ll k$ with $a,b\neq 2$, $a+b\geq 10$, and $T$ a graph on $[a]\cup [b]$For any $a,b\ll k$ with $a,b\neq 2$, $a+b\geq 10$, we can pick disjoint zero-sum subsets $A,B\subseteq \mathbb{Z}_2^k$ with $|A|=a, |B|=b$ such that  $x+y\neq x'+y'$ for any subsets $\{x,y\},\{x',y'\}\subseteq A\cup B$ unless $\{x,y\}=\{x',y'\}$.
\end{lemma}

\begin{theorem}\label{Theorem_harmonious_tree_characterization}
Let $T$ be a tree with $\Delta(T)\leq \Delta$  and $G$ an abelian group. There is a rainbow copy of $T$ in $K_G$ if, and only if, we have none of the following:
\begin{enumerate}[(1)]
\item $G=\mathbb{Z}_2^m$, $m\geq 2$  and  $T$ is a path (or equivalently $T$ has precisely two vertices of odd degree).
\item $G=\mathbb{Z}_2^m$, $m\geq 2$  and  $T$ has precisely two vertices of even degree.
\item $G=\mathbb{Z}_{m_1}\times\dots \times \mathbb{Z}_{m_k}$ and $V(T)=\{v_1, \dots, v_n\}$, with $v_1v_2\in E(T)$ and  $d_T(v_1), d_T(v_2)\equiv 0\pmod {m_i}$, $d_T(v_3), \dots, d_T(v_n)\equiv 1\pmod {m_i}$ for all $i$.
\item $G=\mathbb{Z}_2^m$, $m\geq 2$ , $T$ contains precisely $4$ vertices of  even degree and has a perfect matching when restricted to these $4$ vertices.
\end{enumerate}
\end{theorem}

\begin{proof}

``Only if'' direction:\\
Suppose that $G=\mathbb{Z}_2^m$. Let $V_{odd}, V_{even}$ be the sets of odd/even degree vertices of $T$.
Suppose that there is a rainbow embedding $\phi$ of $T$ into $V(G)$. Then we must have $C(\phi(T))=G\setminus \{0\}$ since the colour $0$ doesn't appear on any edges of $K_{\mathbb{Z}_2^m}$.
Then $\sum_{v\in V(T)}\phi(v)=\sum G=0$ (using that $m\geq 2$) and $\sum_{v\in V(T)}d_T(v)\phi(v)=\sum C(\phi(T))=\sum G-0=0$. Since $\sum_{v\in V_{odd}}\phi(v)=\sum_{v\in V_{odd}}d_T(v)\phi(v)=\sum_{v\in V(T)}d_T(v)\phi(v)$ and $\sum_{v\in V_{even}}\phi(v)= \sum_{v\in V(T)}\phi(v)-\sum_{v\in V_{odd}}\phi(v)$, we get that $\sum_{v\in V_{odd}}\phi(v), \sum_{v\in V_{even}}\phi(v)=0$. Since all the vertices of $\phi(V_{even}), \phi(V_{odd})$ must be distinct, this means that $|V_{odd}|, |V_{even}|\neq 2$ (because in $\mathbb{Z}_2^m$ we cannot have two distinct elements adding to $0$). If $|V_{even}|=4$, then we get that $T[V_{even}] $ doesn't have a perfect matching, since otherwise the two edges of this matching must have the same colour in the embedding.


Suppose that ``$G=\mathbb{Z}_{m_1}\times\dots \times \mathbb{Z}_{m_k}$ and $V(T)=\{v_1, \dots, v_n\}$, where $d_T(v_3), \dots, d_T(v_n)\equiv 1\pmod {m_i}$ for all $i$, and $v_1v_2\in E(T)$ and $d_T(v_1), d_T(v_2)\equiv 0\pmod {m_i}$ for all $i$''.
Suppose for contradiction that there is a rainbow embedding $\phi$ of $T$ into $V(G)$. Let $c$ be the unused colour. We have $\sum G=\sum_{v\in V(T)} \phi(v)$ and $\sum G-c=\sum_{v\in V(T)} d_T(v)\phi(v)$. Subtracting gives $c=\sum_{v\in V(T)} (1-d_T(v))\phi(v)$. Since $d_T(v_3), \dots, d_T(v_n)\equiv 1\pmod {m_i}$ and $d_T(v_1), d_T(v_2)\equiv 0\pmod {m_i}$ for all $m_i$ we have that $(1-d_T(v_i))\phi(v_i)=0$ in $G$ for $i=3, \dots, m$ and $(1-d_T(v_1))\phi(v_1)=\phi(v_1)$, $(1-d_T(v_2))\phi(v_2)=\phi(v_2)$.
 This gives  $c=\sum_{v\in V(T)} (1-d_T(v))\phi(v)=(1-d_T(v_1))\phi(v_1)+(1-d_T(v_2))\phi(v_2)=\phi(v_1)+\phi(v_2)$. But $\phi(v_1)+\phi(v_2)$ is also the colour of the edge $v_1v_2$ contradicting that $c$ is not used on $\phi(T)$.




``If'' direction:\\
Suppose $G\neq \mathbb{Z}_2^m$. If $T$ is a path, let $T_{core}$  be an independent set consisting  of both leaves, and $6$ vertices of degree $2$ (noting that this is a core of $T$). Otherwise, let  $T_{core}$ be a core of $T$ of size $\leq 12\Delta$ given by Observation~\ref{Observation_small_core}, noting that $T_{core}$ will then contain $\geq 3$ leaves of $T$ (since $T$ is not a path it has $\geq 3$ leaves. Now depending on whether I or II occurs, $T_{core}$ contains either all the leaves or at least $6$  leaves). Label $V(T_{core})=\{v_1, \dots, v_k, w_1, \dots, w_t, u_1, u_2, \dots, u_s\}$ where
for all $i$, $d(v_i)\not\equiv 1 \pmod G$,  $d(w_i)\equiv 1 \pmod G$ with $d(w_i)\neq 1$, and $u_1, \dots, u_s$ are leaves.  
Let $T_{core}^{v}=T[\{v_{1}, \dots, v_k\}]$ and
$T_{core}^{w}=T[\{v_{1}, \dots, v_k, w_1, \dots, w_t\}]$
%Define  $c_{special}:=(-d(v_1)+1)x_1+\dots+(-d(v_k)+1)x_k$.

\begin{claim}\label{Claim_harmonious_tree_characterization1}
There is a  $\psi:\{v_1, \dots, v_k\}\to V(K_G)$ which is a rainbow embedding of $T_{core}^v$ not using the colour $c_{special}:=(-d(v_1)+1)\psi(v_1)+\dots+(-d(v_k)+1)\psi(v_k)$.
\end{claim}
\begin{proof}
Note that if  $k=2$, then $v_1v_2\not\in E(T)$ --- otherwise the degrees of all vertices in $T_{core}$ other than $v_1,v_2$ must be $\equiv 1 \pmod G$. This would imply that the degrees of all vertices in $T$ other than $v_1,v_2$ are $\equiv 1 \pmod G$ (since $T_{core}$ has a representative vertex of every degree occurring in $T$ by the definition of ``core''), and hence (3) would hold.

Thus we can apply Lemma~\ref{Lemma_kastv_simple_tree} we get some simple $x\in \mathbb{G}^k$ with  $x\ast_T d$ simple. Define $\psi$ to embed $v_i$ to $x_i$ for all $i$. Since $x$ is simple, this is an injection. 
The multiset of colours it uses is $\{x_i+x_j: v_iv_j\in E(T)\}=x\ast_T d\setminus\{c_{special}\}$. Since $x\ast_T d$ is simple, we get that these colours are all distinct from each other and from $c_{special}$.
\end{proof}


\begin{claim}\label{Claim_harmonious_tree_characterization2}
We can extend $\phi$ to  $\theta:\{v_1, \dots, v_k, w_1, \dots, w_t\}\to V(K_G)$ which is a rainbow embedding of $T_{core}^w$ not using the colour $c_{special}$.
\end{claim}
\begin{proof}
For $i=0, \dots, t$ set $V_i=\{v_1, \dots, v_k, w_1, \dots, w_i\}$. 
We build rainbow embeddings $$\theta_i:\{v_1, \dots, v_k, w_1, \dots, w_i\}\to V(K_G)$$ one by one. 
Start with $\theta_0=\psi$. To build $\theta_{i+1}$ from $\theta_{i}$:
Pick $\theta_{i+1}(w_{i+1})$ to be anything outside $\theta_{i}(V_{i})\cup (\theta_{i}(V_{i})+\theta_{i}(V_{i})-\theta_{i}(V_{i}))\cup (c_{special}-\theta_{i}(V_{i})$ (there's space to do this since $|\theta_{i}(V_{i})\cup (\theta_{i}(V_{i})+\theta_{i}(V_{i})-\theta_{i}(V_{i}))\cup (c_{special}-\theta_{i}(V_{i}))|\leq (k+t)+(k+t)^3+(k+t)\leq 3|V(T_{core})|^3\ll n$). 
Note that $\theta_{i+1}$ is an injection since $\theta_{i}$ was one, and $\theta_{i+1}(w_{i+1})\not\in \theta_{i}(V_{i})$. Also $\theta_{i+1}$ is a rainbow embedding since $\theta_{i}$ was one, and the new colours used by $\theta_{i+1}$ are contained in $\theta_{i+1}(w_{i+1})+\theta_{i}(V_i)$ which is disjoint from $C(\theta_i(T[V_i]))\subseteq \theta_{i}(V_i)+\theta_{i}(V_i)$ (due to $\theta_{i+1}(w_{i+1})\not\in (\theta_{i}(V_i)+\theta_{i}(V_i)-\theta_{i}(V_i))$). 
Finally, $\theta_{i+1}$ doesn't use the colour $c_{special}$ since $\theta_{i+1}(w_{i+1})+\theta_{i}(V_i)$ is disjoint from $\{c_{special}\}$ (this is equivalent to $\theta_{i+1}(w_{i+1})\not\in (c_{special}-\theta_{i}(V_i))$).
\end{proof}

\begin{claim}\label{Claim_harmonious_tree_characterization3}
We can extend $\theta$ to  $\phi:\{v_1, \dots, v_k, w_1, \dots, w_t, u_1, \dots, u_s\}\to V(K_G)$ which is a rainbow embedding of $T_{core}$ not using the colour $c_{special}$ and satisfying $\sum V(\phi(T_{core}))=\sum G$.
\end{claim}
\begin{proof}
Recall $u_1, \dots, u_s$ are leaves with $s\geq 2$. Let $N:=\bigcup_{i=1}^sN_{T}(u_i)\cap V(T_{core})$, noting that when $s=2$, we have ensured $N=\emptyset$. Let $g:=\sum G- \sum im_{\theta}$, $F_1= im_{\theta}\cup (im_{\theta}+im_{\theta}-im_{\theta})\cup (c_{special} -im_{\theta})$ and $F_2=N-N$, noting that $|F_1|, |F_2|\leq 3|T_{core}|^3\leq 3(12\Delta)^3\ll n$ and that when $s=2$ we have $F_{2}=\{0\}$. 
Depending on whether $s=2$ or not, use Lemma~\ref{Lemma_y1y2_separate} or~\ref{Lemma_y1y2...ys_separate} to pick  $y_1, \dots, y_s\not\in F_1$ with $y_1+\dots+y_s=g$ and $y_i-y_j\not\in F_2$ for $i\neq j$. 
Define $\phi$ to agree with $\theta$ on $\{v_1, \dots, v_k, w_1, \dots, w_t\}$ and to have $\phi(u_i) = y_i$. 
This is an injection because $y_1, \dots, y_s$ are distinct and outside $im_{\theta}$. When $s=2$, there are no edges in $T_{core}$ touching $u_1, \dots, u_s$, so we have a rainbow embedding in that case.  When $s\geq 3$, the  colours of new edges used by $\phi$ (i.e. the colours edges $\phi(xy)$ with $xy\not\in T_{core}^w$) are contained in $\{im_{\theta}+y_i: i=1, \dots, s\}$ (here, we're using that $u_1, \dots, u_s$ is an independent set due to it being a set of leaves of a tree $T$). These colours are all distinct from each other (since $y_i-y_j\not\in im_{\theta}-im_{\theta}$), from the colours of $\theta(T_{core}^w)$ (since $y_i\not \in im_{\theta}+im_{\theta}-im_{\theta}$), and from $c_{special}$ (since  $y_i\not \in c_{special}-im_{\theta}$).
Finally we have $\sum_{v\in V(T_{core})}\phi(v)
=\sum im_{\theta} +\sum_{i=1}^s y_i
=\sum G$
\end{proof}
Set $C_{target}=C(K_G)\setminus\{c_{special}\}$ and $V_{target}=V(K_G)$, noting that $\sum C_{target}=\sum G-c_{special}$ and $\sum V_{target}=\sum G=\sum V(\phi(T_{core}))$. Using that vertices  $v\in V(T_{core})\setminus\{v_1, \dots, v_k\}$ have $d_T(v)\equiv 1 \pmod G$, we get
\begin{align*}
\sum_{v\in V(T_{core})}d_T(v)\phi(v)&= \sum_{v\in V(T_{core})}(d_T(v)-1)\phi(v)+\sum_{v\in V(T_{core})}\phi(v)= \sum_{v\in V(T_{core})}(d_T(v)-1)\phi(v)+\sum G\\
&=\sum_{i=1}^{k}(d_T(v_k)-1)\phi(v_k)+\sum G
=\sum G-c_{special}=\sum C_{target}
\end{align*}
 Thus the embedding $\phi$ satisfies Theorem~\ref{Theorem_embed_tree_prescribed_sets} (ii), and hence we get a rainbow embedding of $T$ into $(V_{target}, C_{target})$ (and hence into $K_G$). 


Suppose $G= \mathbb{Z}_2^m$. Use Observation~\ref{Observation_small_core} to get a core $T_{core}$ of $T$ of order $12\Delta$. Set $C_{target}=C(K_G)\setminus\{0\}, V_{target}=V(K_G)$, noting that these both have zero sum (since $\sum \mathbb{Z}_2^m=0$ for $m\geq 2$).
Let $A=\{v_1, \dots, v_a\}$ be the odd degree vertices in $T_{core}$ and $B=\{u_1, \dots, u_b\}$ the even degree vertices. Note that $|A|\neq 2$, as otherwise the degrees $d(v_1)$ and $d(v_2)$ must be exhausted in $T_{core}$ --- since non-exhausted degrees $d$ have $\geq 3$ degree $d$ vertices in every core and thus $v_1, v_2$ are the only odd degree vertices in $T$, and hence $T$ is a path, contradicting (1) not holding. Similarly $|B|\neq 2$ --- since otherwise $u_1$ and $u_2$ would be the only even degree vertices in $T$, contradicting (2) not holding.
 Also note that $|A|+|B|=|V(T_{core})|\geq 10$. If $|A|=4$, note that $T[A]$ can't have  a perfect matching (since leaves can't be connected in a $\ge 3$-vertex tree, for $T[A]$ to have a perfect matching, $A$ must have $\le 2$ leaves. But the only tree with $\le 2$ leaves is a path which doesn't have $4$ odd degree vertices). If $|B|=4$, note that $T[B]$ can't have  a perfect matching, as otherwise we'd have (4).
 
 Use Lemma~\ref{Lemma_zero_sum_sets_Z2k} to get a rainbow embedding $\phi$ of $T[A\cup B]$ with $\sum \phi(A)=\sum \phi(B)=0$. This ensures that $\sum_{v\in V(T_{core})}d_T(v)\phi(v), \sum \phi(V(T_{core}))=0=\sum V_{target}=\sum_{C_{target}}$ and hence by Theorem~\ref{Theorem_embed_tree_prescribed_sets},  we get a rainbow embedding of $T$ in $K_G$.
\end{proof}

\end{document}
